\section{Estrategia de pruebas}
\label{sec:pruebas}

\subsection{Objetivo y organización}
La validación persigue tres metas: (i) comprobar que las tres colas cumplen el mismo contrato; (ii) verificar que Dijkstra devuelve distancias óptimas y rutas reconstruibles; y (iii) asegurar que la carga de datos, la CLI, el servicio meteorológico y la aplicación funcionan sin red y rechazan entradas inválidas. La suite se ejecuta con \texttt{pytest} y se organiza por módulo, no por integrante, para que cualquier cambio pueda evaluarse de forma local.

\begin{lstlisting}[language=bash, caption={Ejecución de la suite completa}]
python -m pytest python/tests -q
\end{lstlisting}

La Tabla~\ref{tab:suite_pruebas} resume las 278 pruebas de la suite, incluida la validación de proveedores simulados, los escenarios de ruta y las cinco secciones de la web. El conteo incluye casos parametrizados: una misma función de prueba puede ejecutarse sobre varias colas, tamaños o entradas inválidas. En el entorno de revisión, \texttt{278 passed}; la duración depende de la máquina.

\subsection{Correspondencia con los contenidos requeridos}
La estrategia cubre explícitamente las cinco categorías solicitadas para el informe final. La Tabla~\ref{tab:categorias_prueba} relaciona cada categoría con su evidencia implementada y con los casos que se verifican.

\begin{table}[H]
    \centering
    \fontsize{9}{11}\selectfont
    \renewcommand{\arraystretch}{1.15}
    \caption{Cobertura de las categorías de prueba requeridas}
    \label{tab:categorias_prueba}
    \begin{tabularx}{\textwidth}{@{}>{\bfseries\raggedright\arraybackslash}p{3.0cm} >{\raggedright\arraybackslash}p{4.0cm} X@{}}
        \toprule
        \textbf{Categoría} & \textbf{Evidencia principal} & \textbf{Casos cubiertos} \\
        \midrule
        Pruebas unitarias & \texttt{test\_priority\_queues.py}, \texttt{test\_dijkstra.py} & Operaciones de las colas, relajación de aristas, extracción mínima, claves repetidas, cola vacía y distancias básicas. \\
        Pruebas de integración & \texttt{test\_cli.py}, \texttt{test\_osm\_loader.py}, \texttt{test\_weather.py}, \texttt{test\_streamlit\_app.py} & Integración de CLI, carga de snapshots OSM, servicio meteorológico con caché y vistas principales de Streamlit sin depender de Internet. \\
        Validación & \texttt{test\_equivalence.py}, \texttt{test\_validation\_io.py} & Comparación con Bellman-Ford, Dijkstra con arreglo, NetworkX y certificado de optimalidad con rutas reconstruibles. \\
        Casos límite & \texttt{test\_dijkstra\_manual.py}, \texttt{test\_dijkstra.py} & Un solo nodo, origen sin salidas, nodos inalcanzables, lazos, aristas paralelas, pesos cero y valores extremos. \\
        Casos adversos & \texttt{test\_validation\_io.py}, \texttt{test\_priority\_queues.py}, \texttt{test\_a\_star.py} & Pesos negativos, ciclos negativos, archivos corruptos, estructuras incompletas, claves inválidas, ciclos de predecesores y heurísticas no admisibles. \\
        \bottomrule
    \end{tabularx}
\end{table}

\begin{table}[H]
    \centering
    \fontsize{9}{11}\selectfont
    \renewcommand{\arraystretch}{1.12}
    \caption{Distribución de la suite automatizada}
    \label{tab:suite_pruebas}
    \begin{tabularx}{\textwidth}{@{}>{\raggedright\arraybackslash}p{5.3cm} r >{\raggedright\arraybackslash}X@{}}
        \toprule
        \textbf{Archivo} & \textbf{Pruebas} & \textbf{Propósito principal} \\
        \midrule
        \texttt{test\_priority\_queues.py} & 61 & Contrato común, orden de extracción, claves inválidas, cortes en cascada y comparación con \texttt{heapq}. \\
        \texttt{test\_equivalence.py} & 48 & Equivalencia con Bellman-Ford, Dijkstra con arreglo y NetworkX; empates y escalas de 1000 y 3000 nodos. \\
        \texttt{test\_validation\_io.py} & 45 & Validación de JSON, entradas adversas y certificado, incluidos ciclos de predecesores de peso cero. \\
        \texttt{test\_dijkstra\_manual.py} & 36 & Traza manual de Cusco, CLRS 24.6, nodos inalcanzables, lazos, aristas paralelas y pesos extremos. \\
        \texttt{test\_cli.py} & 25 & Comandos, códigos de salida, generación, verificación, exportación y modo interactivo. \\
        \texttt{test\_dijkstra.py} & 11 & Casos básicos, límites y equivalencia a escala. \\
        \texttt{test\_a\_star.py} & 8 & Optimalidad, consistencia, ausencia de coordenadas y regresión de heurística no admisible. \\
        \texttt{test\_benchmark\_protocol.py} & 4 & Estadísticas robustas, perfiles y esquema del benchmark canónico. \\
        \texttt{test\_osm\_loader.py} & 4 & Conversión determinista, métricas, caché y rechazo de snapshots inválidos. \\
        \texttt{test\_weather.py} & 12 & Multiplicador $\Omega_{\text{clima}}$, clima adverso, respuestas simuladas de Open-Meteo y OpenWeatherMap, presión de suelo, hora faltante, manejo de 401 y modo desconectado. \\
        \texttt{test\_traffic.py} & 10 & Escenarios horarios, sensibilidad por tipo de vía, sinergia clima-tráfico y respuestas TomTom simuladas sin inventar velocidades. \\
        \texttt{test\_applied\_analysis.py} & 3 & Nodo más cercano, escalamiento uniforme y rechazo de factores inválidos. \\
        \texttt{test\_streamlit\_app.py} & 11 & Figuras con y sin teselas, costos, escenario histórico, validación de datos, cinco secciones, contraste con $A^*$ y rechazo de clave implícita. \\
        \midrule
        \textbf{Total} & \textbf{278} & Suite completa de pruebas automatizadas. \\
        \bottomrule
    \end{tabularx}
\end{table}

\subsection{Patrón Arrange-Act-Assert}
Cada prueba sigue tres fases explícitas:
\begin{itemize}
    \item \textbf{Arrange.} Se construye la instancia, la cola o el archivo temporal, y se declara el resultado esperado.
    \item \textbf{Act.} Se invoca la operación bajo prueba: \texttt{insert}, \texttt{extract\_min}, \texttt{decrease\_key}, \texttt{dijkstra\_solver} o un comando de la CLI.
    \item \textbf{Assert.} Se comprueba igualdad exacta o aproximada, orden, excepción esperada o código de salida.
\end{itemize}
Las comparaciones de distancias usan una tolerancia relativa de $10^{-6}$ para absorber el redondeo de punto flotante; las comparaciones de rutas verifican el costo, no la identidad de la secuencia, porque puede haber varios caminos óptimos.

\subsection{Pruebas del contrato de colas}
Las pruebas de \texttt{test\_priority\_queues.py} se parametrizan sobre las tres estructuras y cubren:
\begin{enumerate}
    \item extracción en orden creciente y vaciado completo;
    \item \texttt{decrease\_key} hacia el frente, desde $\infty$ y sobre un elemento ausente;
    \item rechazo de inserciones duplicadas y de claves NaN, cadenas, \texttt{None} o booleanos;
    \item extracción sobre cola vacía, elemento único y claves repetidas;
    \item comparación aleatoria contra una implementación de referencia, con seis variantes de heap y tres semillas;
    \item ejercitación específica de cortes en cascada y consolidación del Fibonacci heap.
\end{enumerate}
La comparación con \texttt{heapq} se usa solo como oráculo de orden, no como estructura del solver. Esto separa la corrección del contrato de la decisión de diseño de implementar las colas desde cero.

\subsection{Oráculos de caminos mínimos}
La corrección se contrasta con tres referencias independientes:
\begin{itemize}
    \item \textbf{Bellman-Ford.} Tolerante a pesos negativos y capaz de detectar ciclos negativos; se usa en 40 grafos aleatorios de $5$ a $122$ nodos, con y sin pesos enteros, ceros y densidad $4|V|$.
    \item \textbf{Dijkstra con arreglo no ordenado.} Referencia $O(|V|^2+|E|)$ sobre una red de 300 nodos; comprueba que la cola no altera las distancias.
    \item \textbf{NetworkX 3.7.} Oráculo externo sobre una cuadrícula dirigida de $8\times 8$; evita que un error compartido por las tres implementaciones pase inadvertido.
\end{itemize}
Además, \texttt{test\_dijkstra\_manual.py} compara contra valores calculados a mano en la instancia de seis nodos y en la Figura~24.6 de \textcite{cormen2022introduction}. Esa prueba detecta errores de relajación que un oráculo automático podría reproducir si compartiera el mismo sesgo de implementación.

\subsection{Certificado de optimalidad}
Para pesos no negativos, un vector $d$ es óptimo si y solo si cumple cuatro condiciones:
\begin{enumerate}
    \item $d[s]=0$;
    \item para toda arista $(u,v)$ con $d[u]$ finito, $d[v]\le d[u]+w(u,v)$;
    \item los predecesores de cada vértice con distancia finita forman una cadena de aristas ajustadas que llega a $s$; no basta con que cada arista sea ajustada si se forma un ciclo de peso cero desconectado;
    \item cada camino reportado empieza en $s$, termina en $v$, usa aristas existentes y cuesta $d[v]$; los inalcanzables tienen $d[v]=\infty$ y camino vacío.
\end{enumerate}
La comprobación de desigualdades, aristas ajustadas y conexión de predecesores al origen recorre vértices y aristas en $O(|V|+|E|)$ sin volver a ejecutar Dijkstra. Si se entregan caminos materializados, leer y comprobar cada secuencia añade $O(L)$, donde $L$ es la suma de sus longitudes; no se atribuye esa parte a $O(|V|+|E|)$. El benchmark y la consulta OSM verifican la salida fuera de la región cronometrada. \texttt{test\_validation\_io.py} prueba salidas alteradas, ciclos de peso cero con distancia falsa, caminos incompletos y cadenas largas sin recursión.

\subsection{Casos límite}
La suite cubre los siguientes casos límite:
\begin{itemize}
    \item grafo de un solo nodo, origen sin aristas salientes y nodos inalcanzables;
    \item lazos, aristas paralelas, pesos cero, pesos muy grandes y muy pequeños, e identificadores enteros;
    \item cadenas largas sin recursión, para verificar que la reconstrucción de rutas es iterativa;
\end{itemize}

\subsection{Casos adversos y regresión}
La suite cubre los siguientes casos adversos:
\begin{itemize}
    \item ciclos negativos, que deben ser rechazados por Dijkstra y detectados por Bellman-Ford;
    \item archivos inexistentes, extensión incorrecta, JSON mal formado, codificación no UTF-8 y estructuras incompletas.
\end{itemize}
La prueba de regresión de $A^*$ merece mención aparte. Una versión anterior aplicaba coordenadas de Cusco a cualquier grafo y producía una heurística no admisible. El caso de regresión construye un grafo con coordenadas que violan la cota y comprueba que la implementación actual rechaza la heurística o cae a $h=0$; no se limita a comparar el resultado final.

\subsection{Pruebas de integración sin red}
Las pruebas de OSM, clima y tráfico no dependen de Internet. \texttt{test\_osm\_loader.py} usa un \textit{snapshot} temporal, comprueba conversión determinista y caché. \texttt{test\_weather.py} simula respuestas y comprueba respaldo local y distinción entre presión de suelo y del nivel del mar. \texttt{test\_traffic.py} rechaza datos de velocidad faltantes o inválidos. Once pruebas de la aplicación cubren resultados y figuras locales, costos equivalentes, escenarios uniformes, carga de benchmark, acceso sin red a cinco secciones, contraste de cuatro ejecuciones, observación histórica y rechazo de una clave implícita mediante \texttt{AppTest}.

\subsection{Criterio de aceptación}
Una versión se considera aceptable si: (i) las 278 pruebas pasan; (ii) el certificado de optimalidad se cumple en todas las instancias del benchmark; (iii) las tres colas producen las mismas distancias; y (iv) la aplicación y la CLI funcionan con datos locales aunque no haya red. No se reporta cobertura porque no se midió.
